\documentclass[a4paper, serif, 10pt]{moderncv}

%% moderncv
% style and color
\moderncvstyle{banking}
\moderncvcolor{black}

%% page layout/margins
\usepackage[top=2.5cm, bottom=2.5cm, left=1.5cm, right=1.5cm]{geometry}

\nopagenumbers{}

%% Babel config for Brazilian Portuguese language
% ref:
% - https://latex3.github.io/babel/guides/locale-portuguese.html
\usepackage[brazilian]{babel}

%% fontspec
\usepackage{fontspec}

\IfFontExistsTF{Linux Libertine}{
  \setmainfont[Ligatures={TeX, Common}, Numbers=OldStyle]{Linux Libertine}
}{}
\IfFontExistsTF{Linux Libertine O}{
  \setmainfont[Ligatures={TeX, Common}, Numbers=OldStyle]{Linux Libertine O}
}{}

%% CTeX
% CJK support, used to show CJK characters in the resume
%
% - fontset=none: disable builtin fontset but instead set the CJK font manually
% - heading=false: disable ctex heading
% - punct=kaiming: use kaiming punctuations styles for CJK
% - scheme=plain: use plain scheme, do not override \normalsize font size
% - space=auto: space settings for CJK characters
%
% ref:
% - http://ctan.mirrorcatalogs.com/language/chinese/ctex/ctex.pdf
\usepackage[UTF8, heading=false, punct=kaiming, scheme=plain, space=auto]{ctex}

\IfFontExistsTF{Noto Serif CJK SC}{
  \setCJKmainfont{Noto Serif CJK SC}
}{}
\IfFontExistsTF{Noto Sans CJK SC}{
  \setCJKsansfont{Noto Sans CJK SC}
}{}

%% line spacing
\usepackage{setspace}
\setstretch{1.125}

%% URL styling - use normal text instead of monospace
\urlstyle{same}

% Auto-underline all links
% Defer until \begin{document} so hyperref is already loaded
\AtBeginDocument{%
  \let\oldhref\href
  \renewcommand{\href}[2]{\oldhref{#1}{\underline{#2}}}%
}

%% Patch \httplink and \httpslink to handle full URLs
% The original moderncv macros always prepend http:// or https:// to URLs.
% This causes issues when the URL already has a protocol (e.g., https://example.com)
% resulting in malformed URLs like https://https://example.com.
% This patch checks if the URL already contains "://" and uses it directly if so.
% Note: We use \str_set:Nx (x-expansion) to expand macros like \@homepage first.
\ExplSyntaxOn
\str_new:N \g_yamlresume_url_str

\RenewDocumentCommand{\httpslink}{O{}m}{
  \str_set:Nx \g_yamlresume_url_str {#2}
  \str_if_in:NnTF \g_yamlresume_url_str {://}
    { \url{#2} }
    { \url{https://#2} }
}

\RenewDocumentCommand{\httplink}{O{}m}{
  \str_set:Nx \g_yamlresume_url_str {#2}
  \str_if_in:NnTF \g_yamlresume_url_str {://}
    { \url{#2} }
    { \url{http://#2} }
}
\ExplSyntaxOff

\name{Beatriz Carvalho}{}
\title{Analista de Dados | SQL, Python e Power BI}
\phone[mobile]{+55 11 98888-7766}
\email{beatriz.carvalho.dados@gmail.com}
\homepage{https://www.beatrizcarvalho.dev.br}

\address{São Paulo, SP, Brasil, 01310-100}{}{}

\extrainfo{{\small \faLinkedin}\ \href{https://www.linkedin.com/in/beatriz-carvalho-dados}{@beatriz-carvalho-dados} {} {} {} • {} {} {}
{\small \faGithub}\ \href{https://github.com/beatrizcarvalho}{@beatrizcarvalho}}

\begin{document}

\maketitle

\section{Informações Básicas}

\cvline{}{Analista de Dados com mais de 5 anos de experiência nos setores de varejo,
finanças e tecnologia. Atuo na transformação de dados brutos em insights
acionáveis, apoiando decisões estratégicas com análises descritivas,
dashboards executivos e automação de relatórios.
%
\begin{itemize}
\item Especialista em SQL, Python, Power BI e estatística descritiva.
\item Experiência em ETL, modelagem dimensional e storytelling com dados.
\item Histórico de redução de 30\% no tempo de elaboração de relatórios e
aumento de 12\% na conversão de campanhas analisadas.
\item Perfil colaborativo, com forte comunicação entre áreas de negócio e
tecnologia.
\end{itemize}}
\section{Formação}

\cventry{fev. de 2015–dez. de 2018}
        {Bacharelado, Estatística e Ciência de Dados, Nota: 8.7}
        {Universidade de São Paulo}
        {\url{https://www.ime.usp.br}}
        {}
        {\begin{itemize}
\item Formação sólida em estatística, probabilidade e análise de dados.
\item Projetos acadêmicos com modelagem preditiva, testes de hipóteses e
visualização de dados.
\item Iniciação científica em análise de séries temporais aplicada a dados
econômicos.
\end{itemize}
\textbf{Cursos}: Estatística Descritiva, Probabilidade I e II, Inferência Estatística, Modelos Lineares, Introdução à Programação, Banco de Dados, Aprendizado de Máquina, Séries Temporais}
\section{Experiência Profissional}

\cventry{jan. de 2023–Atual}
        {Analista de Dados Pleno}
        {TechRetail Brasil}
        {\url{https://www.techretail.com.br}}
        {}
        {\begin{itemize}
\item Lidero a construção de dashboards executivos em Power BI para as áreas
de vendas, marketing e operações.
\item Desenvolvo pipelines de ETL em Python e SQL integrando dados de ERP,
CRM e Google Analytics.
\item Realizo análises de desempenho de campanhas e comportamento de
clientes, com recomendações que aumentaram a conversão em 12\% no
último ano.
\item Automatizo relatórios semanais e mensais, reduzindo em 30\% o tempo de
elaboração.
\item Atuo como ponte entre as áreas de negócio e tecnologia para definição
de métricas e indicadores.
\end{itemize}
\textbf{Palavras-chave}: Power BI, SQL, Python, ETL, Storytelling}

\cventry{mar. de 2021–dez. de 2022}
        {Analista de Dados Júnior}
        {CrediZap Financeira}
        {\url{https://www.credizap.com.br}}
        {}
        {\begin{itemize}
\item Criei relatórios gerenciais de crédito e inadimplência no Tableau e
Excel.
\item Modelei dados da carteira de clientes em SQL Server para subsidiar
análises de risco.
\item Automatizei extrações e tratamento de dados com Python, reduzindo
retrabalho operacional.
\item Colaborei com cientistas de dados em projetos de score de crédito e
segmentação de clientes.
\end{itemize}
\textbf{Palavras-chave}: Tableau, SQL Server, Excel, Análise de Crédito, Python}

\cventry{jun. de 2019–fev. de 2021}
        {Estagiária de Inteligência de Mercado}
        {DataSeg Consultoria}
        {\url{https://www.dataseg.com.br}}
        {}
        {\begin{itemize}
\item Suportei consultores em projetos de pesquisa de mercado e análise de
concorrência.
\item Estruturei bases de dados em Excel e Access para consolidação de
indicadores.
\item Produzi materiais de apoio com análises descritivas e gráficos para
apresentações a clientes.
\end{itemize}
\textbf{Palavras-chave}: Pesquisa de Mercado, Excel, Análise Descritiva, Relatórios}
\section{Idiomas}

\cvline{Português}{Nativo ou Bilíngue \hfill \textbf{Palavras-chave}: Português do Brasil}
\cvline{Inglês}{Proficiência Profissional Fluente \hfill \textbf{Palavras-chave}: TOEFL ITP 587, Leitura e conversação técnica}
\cvline{Espanhol}{Proficiência Profissional Limitada \hfill \textbf{Palavras-chave}: Espanhol intermediário}
\section{Competências}

\cvline{Análise de Dados}{Especialista \hfill \textbf{Palavras-chave}: SQL, Python, Pandas, Estatística, Excel}
\cvline{Visualização de Dados}{Avançado \hfill \textbf{Palavras-chave}: Power BI, Tableau, DAX, Looker Studio, Storytelling com Dados}
\cvline{Bancos de Dados}{Avançado \hfill \textbf{Palavras-chave}: PostgreSQL, MySQL, SQL Server, BigQuery, NoSQL}
\cvline{Estatística e Machine Learning}{Intermediário \hfill \textbf{Palavras-chave}: Regressão, Testes A/B, Séries Temporais, scikit-learn}
\section{Prêmios}

\cventry{dez. de 2023}
        {TechRetail Brasil}
        {Prêmio Excelência em Dados}
        {}
        {}
        {Reconhecimento interno pela liderança no projeto de painel de
indicadores estratégicos, que impactou positivamente as decisões da
diretoria de varejo.}
\section{Certificados}

\cventry{mar. de 2023}
        {Microsoft}
        {Certificação Microsoft PL-300 - Power BI Data Analyst}
        {\url{https://learn.microsoft.com/pt-br/credentials/certifications/data-analyst-associate/}}
        {}
        {}

\cventry{jul. de 2021}
        {Google}
        {Google Data Analytics Professional Certificate}
        {\url{https://www.coursera.org/professional-certificates/google-data-analytics}}
        {}
        {}
\section{Projetos}

\cventry{fev. de 2023–mar. de 2023}
        {Desenvolvimento de um dashboard em Power BI para acompanhamento de metas comerciais.}
        {Dashboard de Vendas Interativo}
        {\url{https://github.com/beatrizcarvalho/dashboard-vendas}}
        {}
        {\begin{itemize}
\item Modelei dados de vendas no Power Query e DAX.
\item Criei indicadores de receita, ticket médio e margem por região e
categoria.
\item Publiquei o dashboard como portfólio de análise de dados.
\end{itemize}
\textbf{Palavras-chave}: Power BI, DAX, SQL, Vendas}

\cventry{nov. de 2022–dez. de 2022}
        {Projeto de análise exploratória e predição de churn para uma base de clientes de telecomunicações.}
        {Análise de Churn de Clientes}
        {\url{https://github.com/beatrizcarvalho/churn-analysis}}
        {}
        {\begin{itemize}
\item Apliquei técnicas de análise exploratória e testes de hipóteses.
\item Construí modelo de regressão logística com scikit-learn.
\item Apresentei os principais fatores associados ao cancelamento de clientes.
\end{itemize}
\textbf{Palavras-chave}: Python, Pandas, scikit-learn, Churn}
\section{Interesses}

\cvline{Tecnologia}{Ciência de Dados, Open Source, Visualização de Dados}
\cvline{Esportes}{Corrida, Vôlei, Natação}
\section{Voluntariado}

\cventry{mar. de 2022–dez. de 2023}
        {Mentora de Dados}
        {Programa Meninas na Computação}
        {\url{https://www.meninasnacomputacao.org.br}}
        {}
        {\begin{itemize}
\item Atuei como mentora voluntária em oficinas de introdução à análise de
dados.
\item Auxiliei na elaboração de material didático em Python e Power BI.
\item Incentivei a participação de mulheres em carreiras de tecnologia e
dados.
\end{itemize}}

\end{document}